\chapter{Контроль соответствия требованиям} \label{ch6}

\section{Соответствие требованиям к отчету} \label{ch6:report}

\noindent
\begingroup
\centering
\small
\begin{longtable}[c]{|p{6.2cm}|p{8.8cm}|}
    \caption{Проверка обязательных элементов раздаточного материала}
    \label{tab:requirements-report}
    \\
    \hline
    Требование & Где выполнено \\ \hline
    \endfirsthead
    \hline
    Требование & Где выполнено \\ \hline
    \endhead
    \hline
    Относится к автоматизации или теории тестирования & Тема посвящена автоматизации модульного тестирования JavaScript. \\ \hline
    Принцип IMRAD & Введение, методы, результаты, обсуждение и заключение выделены в структуре отчета. \\ \hline
    Математическая постановка & Раздел~\ref{ch1:math}. \\ \hline
    Иллюстративный пример & Раздел~\ref{ch1:example}. \\ \hline
    Псевдокод в окружении \texttt{algorithm} & Раздел~\ref{ch1:algorithm}, алгоритм~\ref{alg:benchmark}. \\ \hline
    Пошаговое выполнение метода & Раздел~\ref{ch1:steps}. \\ \hline
    Описание разработанного ПО & Разделы~\ref{ch1:software} и~\ref{ch4}. \\ \hline
    Описание средств и методов тестирования & Разделы~\ref{ch1:testing-methods} и~\ref{ch5}. \\ \hline
    Инструкция установки и запуска & Раздел~\ref{ch2:run}. \\ \hline
    Автоматический список источников из BibTeX & Файл \texttt{my\_folder/my\_biblio.bib} и раздел списка источников. \\ \hline
\end{longtable}
\normalsize
\endgroup

\section{Соответствие требованиям к программному обеспечению} \label{ch6:software}

\noindent
\begingroup
\centering
\small
\begin{longtable}[c]{|p{6.2cm}|p{8.8cm}|}
    \caption{Проверка требований к программному обеспечению}
    \label{tab:requirements-software}
    \\
    \hline
    Требование & Реализация \\ \hline
    \endfirsthead
    \hline
    Требование & Реализация \\ \hline
    \endhead
    \hline
    Комментирование кода & Комментарии добавлены в \texttt{src/summary.js} и \texttt{src/benchmark.js} в местах генерации данных и измерения ресурсов. \\ \hline
    Компактность проекта & Основная логика находится в двух файлах: \texttt{summary.js} и \texttt{benchmark.js}. \\ \hline
    Быстрое разворачивание & Установка выполняется командой \texttt{npm install}; проверка --- \texttt{npm run verify}. \\ \hline
    \texttt{requirements.txt} & Файл добавлен в корень \texttt{benchmark\_app}. \\ \hline
    Входные данные переменной величины & Поле \texttt{size} в JSON задает размер массива. \\ \hline
    Импорт CSV или JSON & Реализован импорт JSON. \\ \hline
    Эталонный вход и выход & \texttt{reference\_input.json} и \texttt{reference\_output.json}. \\ \hline
    Авторский пример & \texttt{author\_example.json}. \\ \hline
\end{longtable}
\normalsize
\endgroup

\section{Риски и способы их снижения} \label{ch6:risks}

Основной риск эксперимента связан с нестабильностью ресурсных измерений. Для его снижения рекомендуется выполнять несколько повторов и сравнивать медианные значения. Второй риск --- различие версий фреймворков. Он снижается фиксацией зависимостей в \texttt{package.json}. Третий риск --- различие платформенных средств мониторинга CPU и памяти. Для учебной версии используется переносимая приближенная оценка, а для промышленной версии этот слой можно заменить платформенным монитором.

\section{Итоговая оценка} \label{ch6:summary}

Подготовленный отчет и программное средство покрывают обязательные элементы задания. Работа показывает не только теоретическое сравнение Jest, Vitest и Mocha, но и практический путь получения измерений. Это делает материал полезным для выбора фреймворка в реальном JavaScript-проекте, где важны не только возможности API, но и стоимость регулярного автоматического запуска тестов.
